% SGA 6, Exposición 0: fragmento español alineado con la fuente.
% Alcance: portada y secciones 1--2 completas, antes de la sección 3.
% Fuente francesa: sga6_fr_r1/src/fr.tex, líneas 171--286 inclusive.
% SHA-256 de la fuente francesa completa:
% 0D2158B9FE3B466309A92165727A0B191422E05FBAA0DACD5806E8322C54BA1A.

\section*{Exposición 0}
\begin{center}
\underline{ESBOZO DE UN PROGRAMA PARA UNA TEORÍA DE LAS INTERSECCIONES}\\[3pt]
\underline{SOBRE LOS ESQUEMAS GENERALES}\\[8pt]
\underline{por A. Grothendieck}
\end{center}

La presente exposición es de carácter introductorio, y su lectura no es
lógicamente indispensable para el estudio del Seminario. Se dirige más
especialmente a los lectores familiarizados con la versión provisional del
teorema de Riemann--Roch, tal como se expone en el informe de Borel--Serre
[2] o en el informe de Grothendieck ya mencionado en la Introducción (citado
en adelante como [RRR]), y que se reproduce como apéndice al final de la
presente exposición.

\subsection*{1.}
Recordemos la fórmula de Riemann--Roch para un morfismo propio
\[
f:X\longrightarrow Y
\]
de esquemas cuasiproyectivos y lisos sobre un cuerpo $k$, y para un haz
coherente $F$ sobre $X$, cuya clase en el grupo $K(X)$ de las clases de
haces coherentes sobre $X$ se denota por $\cl(F)$:
\begin{equation}\tag{1.1}
\Todd(T_Y)\,\ch_Y\bigl(f_*(\cl(F))\bigr)=f_*\bigl(\Todd(T_X)\,\ch_X(F)\bigr).
\end{equation}
Esta fórmula es válida en $A(Y)\otimes_{\mathbb Z}\mathbb Q$, donde $A(Y)$
es el anillo de Chow de $Y$, y el $f_*$ del segundo miembro se deduce,
mediante tensorización por $\mathbb Q$, del homomorfismo «~imagen directa de
ciclos~»
\[
f_*:A(X)\longrightarrow A(Y),
\]
mientras que el del primer miembro es la característica de
Euler--Poincaré de $F$ relativa a $f$:
\[
f_*(\cl(F))=\sum_i(-1)^i\cl\bigl(R^if_*(F)\bigr)~;
\]
$\ch_X$, $\ch_Y$ designan los caracteres de Chern sobre $X$, resp. sobre
$Y$, y $T_X$, $T_Y$ los fibrados tangentes a $X$, resp. a $Y$. Como es
sabido, $\Todd(-)$ y $\ch(-)$ son polinomios universales con coeficientes en
$\mathbb Q$ en las clases de Chern del argumento. Como el término constante
de $\Todd(-)$ es $1$, se trata de un elemento invertible para todo valor del
argumento, de modo que la fórmula (1.1), después de multiplicar por
$\Todd(T_Y)^{-1}$, adopta la forma equivalente, más conveniente para nuestro
propósito,
\begin{equation}\tag{1.2}
\ch_Y\bigl(f_*(\cl(F))\bigr)=f_*\bigl(\Todd(T_f)\cdot\ch_X(F)\bigr),
\end{equation}
donde se ha puesto
\begin{equation}\tag{1.3}
T_f=T_X-f^*(T_Y)\in K(X),
\end{equation}
de modo que $T_f$ desempeña el papel de un \underline{fibrado tangente
relativo virtual} de $X$ sobre $Y$. En el caso en que el morfismo $f$ es
liso (es decir, con aplicación tangente sobreyectiva en todo punto), se
tiene simplemente
\[
T_f=T_{X/Y}\qquad\text{(fibrado tangente a lo largo de las fibras)},
\]
mientras que, en el caso en que $f:X\to Y$ es una inmersión, se obtiene
\[
T_f=-\check N_{X/Y},
\]
donde $\check N_{X/Y}$ designa el haz normal de $X$ en $Y$.

Uno de los objetivos principales del Seminario es generalizar (1.2)
simultáneamente en dos direcciones:
\begin{enumerate}[label=\alph*)]
\item Prescindir de la hipótesis de existencia de un cuerpo base $k$.
\item Sustituir las hipótesis de regularidad sobre $Y,X$ por una hipótesis de
«~regularidad local~» de $f$.
\end{enumerate}
Finalmente, a lo largo del camino nos ocuparemos también del problema:
\begin{enumerate}[label=\alph*),start=3]
\item Eliminar las hipótesis de cuasiproyectividad que, en ausencia de un
cuerpo base, se expresan mediante la existencia de módulos invertibles
amplios sobre $X$ y sobre $Y$.
\end{enumerate}

\subsection*{2.}
Examinemos primero la generalización a), conservando, no obstante, las
hipótesis de regularidad y de existencia de módulos invertibles amplios sobre
$X$ y sobre $Y$.

La definición de $K(X)$, $K(Y)$ y del homomorfismo $f_*:K(X)\to K(Y)$ no
presenta entonces ningún problema nuevo, gracias al hecho de que $X$ e $Y$
son regulares. La vía más natural para dar sentido a (1.2) parece consistir,
por tanto, en definir \underline{anillos de Chow} $A(X),A(Y)$ y un
homomorfismo de grupos
\[
f_*:A(X)\longrightarrow A(Y),
\]
en establecer una teoría de clases de Chern que proporcione aplicaciones
\[
c_i:K(X)\longrightarrow A(X),
\]
y análogamente para $Y$, y, finalmente, en describir un elemento que
desempeñe el papel de fibrado tangente relativo virtual
\[
T_f\in K(X).
\]

\subsubsection*{2.1}
En cuanto a este último, se ofrece una definición bastante evidente. Se
utiliza el hecho de que el morfismo $f$, gracias a la hipótesis de existencia
de un módulo invertible amplio sobre $X$ y sobre $Y$, puede factorizarse como
\begin{equation}\tag{2.1}
X\xrightarrow{i}X'\xrightarrow{f'}Y,
\end{equation}
donde $i$ es una inmersión cerrada y $f'$ un morfismo liso; por ejemplo, se
puede tomar $X'=\mathbb P^r_Y$ para un $r$ conveniente. Se pondrá entonces
\begin{equation}\tag{2.2}
\check T_f=\cl({\underline\Omega}^1_{X'/Y})-\cl({\underline N}_{X/X'}),
\end{equation}
donde ${\underline\Omega}^1_{X'/Y}$ es el módulo localmente libre de las
$1$-diferenciales relativas de $X'$ sobre $Y$ (o módulo cotangente relativo
de $X'$ sobre $Y$), y ${\underline N}_{X/X'}$ el módulo conormal de $X$ en
$X'$, igual, por definición, a ${\underline J}/{\underline J}^2$, donde
${\underline J}$ es el ideal sobre $X'$ que define el subesquema cerrado
$X$; como $X$ y $X'$ son regulares, se concluye además que
${\underline N}_{X/X'}$ es también localmente libre. Finalmente,
$\check T_f$ designa el dual de $T_f$. Se comprueba sin dificultad (Exp.
VIII) que el elemento $T_f$ definido por (\sisi{2.1}{2.2}) no depende de la % Canon SGA-SRC-001468
factorización elegida de $f$.

\subsubsection*{2.2}
En cuanto a la definición, bajo las condiciones consideradas, de un anillo
de Chow y de una teoría correspondiente de clases de Chern, aunque es
plausible que tal teoría deba poder desarrollarse siguiendo el modelo de la
ya conocida para los esquemas algebraicos lisos y cuasiproyectivos, parece
que este trabajo no se ha realizado hasta ahora y, a fortiori, no se ha
emprendido en el Seminario. En cambio, se introduce otro anillo, que
desempeñará un papel análogo al del anillo de Chow y cuyo producto tensorial
por $\mathbb Q$ es efectivamente isomorfo a
$A(X)\otimes_{\mathbb Z}\mathbb Q$ cuando $X$ es cuasiproyectivo y liso
sobre un cuerpo $k$. Este anillo es el anillo graduado asociado a $K(X)$
para una filtración adecuada de este anillo.

Una primera filtración natural de $K(X)$, que es la considerada en [RRR], es
la «~\underline{filtración topológica}~» por la codimensión del soporte de
los haces coherentes; $\operatorname{Filt}^{i}_{\mathrm{top}}K(X)$ está
engendrado por los $\cl(F)$ con $F$ coherente y
$\operatorname{codim}(\operatorname{supp}F,X)\ge i$. Una primera cuestión es
la compatibilidad de esta filtración con la estructura de anillo, cuestión
resuelta en [RRR], cuando $X$ es liso sobre un cuerpo $k$, gracias al
«~lema de desplazamiento~» de Chow. Admitido este punto, se obtiene, pues, un
anillo graduado asociado, que aquí denotaremos por
$\Gr^{\bullet}_{\mathrm{top}}(X)$. En [RRR] se define un homomorfismo de
anillos sobreyectivo
\[
\varphi:A(X)\longrightarrow \Gr^{\bullet}_{\mathrm{top}}(X),
\]
(que transforma la clase de un ciclo primo $Z$ en
$\cl(\underline O_Z)$), cuyo núcleo resulta después (Exp. XIV 4.2) ser un
subgrupo de torsión. Además, en [RRR] se muestra cómo las imágenes en
$\Gr^{\bullet}_{\mathrm{top}}(X)$ de las clases de Chern $c_i(F)$ pueden
calcularse directamente en términos de $\cl(F)\in K(X)$ y de las
operaciones $\lambda^i$ de potencia exterior en $K(X)$.

Esto muestra que, para $X$ esquema noetheriano regular provisto de un módulo
invertible amplio, el anillo $\Gr^{\bullet}_{\mathrm{top}}(X)$ es un
sustituto conveniente del hipotético anillo de Chow $A(X)$; tiene además
sobre este último la ventaja de que ciertos cálculos de [RRR] se efectúan
directamente en este anillo y conducen a fórmulas importantes en él, fórmulas
que, incluso para $X$ proyectivo y liso sobre un cuerpo, no están
establecidas actualmente en el propio $A(X)$ (Exp. XIV, 4.3 y 4.4). Sin
embargo, para justificar por completo este punto de vista, habría que
resolver el problema antes señalado de la compatibilidad de la filtración
topológica de $K(X)$ con su estructura de anillo, problema que parece
esencialmente ligado al «~lema de desplazamiento~» de Chow, el cual es a su
vez el ingrediente técnico esencial de la teoría del anillo de Chow que
hemos pretendido evitar.

\subsubsection*{2.3}
Para evitar este problema, se dota a $K(X)$ de una segunda filtración, más
fina que la filtración topológica, que puede describirse en términos
puramente algebraicos a partir de la estructura de $\lambda$-anillo
aumentado de $K(X)$. La definición de esta filtración viene sugerida por la
expresión dada en [RRR] para las clases de Chern de un $x\in K(X)$ en
$\Gr^{\bullet}_{\mathrm{top}}(X)$,
\begin{equation}\tag{2.3}
c_i(x)=\gamma^i(x-\epsilon(x))\mod \operatorname{Filt}^{i+1}_{\mathrm{top}}K(X),
\end{equation}
donde $\gamma^i$ designa cierta combinación lineal de los $\lambda^j$
($j\le i$), explicitada en loc. cit., siendo $\epsilon$ la aumentación. Se
demuestra que $\gamma^i(x-\epsilon(x))$ es efectivamente de filtración
topológica $\ge i$, lo que da sentido a la fórmula (2.3). Se define entonces
la $\lambda$-\underline{filtración} de $K(X)$ mediante los ideales
$\operatorname{Filt}^i K(X)$, donde $\operatorname{Filt}^i K(X)$ está
formado por las sumas de expresiones de la forma
\[
\gamma^{i_1}(x_1-\epsilon(x_1))\cdots \gamma^{i_k}(x_k-\epsilon(x_k)),\qquad i_1+\cdots+i_k\ge i.
\]
El anillo graduado asociado a esta filtración se denotará por
$\Gr^\bullet(X)$, y es este el que sustituirá en el Seminario al anillo de
Chow de $X$. Por otra parte, se puede demostrar que, en el caso aquí
considerado, el homomorfismo canónico
\begin{equation}\tag{2.4}
\Gr^\bullet(X)\longrightarrow \Gr^{\bullet}_{\mathrm{top}}(X)
\end{equation}
es un isomorfismo módulo torsión, lo que justifica el punto de vista según el
cual este anillo, tensorizado por $\mathbb Q$, puede desempeñar en todos los
aspectos el mismo papel que el anillo de Chow tensorizado por $\mathbb Q$.
Por otra parte, se ha hecho exactamente lo necesario para que la fórmula
(2.3) (donde se sustituye $\operatorname{Filt}^i_{\mathrm{top}}$ por
$\operatorname{Filt}^i$) defina una teoría de clases de Chern con valores en
$\Gr^\bullet(X)$, cuyas clases de Chern en
$\Gr^{\bullet}_{\mathrm{top}}(X)$ se deducen mediante el homomorfismo (2.4).
Se demuestra (Exp. V), gracias al hecho (Exp. VI) de que $K(X)$ es un
$\lambda$-anillo \underline{especial} en el sentido de [RRR] (es decir, un
$\lambda$-anillo en el sentido de Exp. V), que esta noción de clase de Chern
satisface todas las propiedades formales habituales, examinadas en [RRR].

\subsubsection*{2.4}
Se observará una ventaja muy apreciable de la definición de
$\Gr^\bullet(X)$ sobre las de $A(X)$ y
$\Gr^{\bullet}_{\mathrm{top}}(X)$: conserva su sentido para un esquema $X$
absolutamente arbitrario y, más generalmente, para un topos localmente
anillado cualquiera $X$, tomando entonces como $K(X)$ el grupo de clases
construido con los módulos localmente libres sobre $X$, que sigue siendo un
$\lambda$-anillo aumentado y, por tanto, está provisto de una filtración
correspondiente que permite definir un graduado asociado $\Gr^\bullet(X)$.
Como precio de esta generalidad, debe señalarse, sin embargo, que, cuando no
se está dispuesto a despreciar los fenómenos de torsión, es decir, a
tensorizar por $\mathbb Q$, las propiedades de $\Gr^\bullet(X)$ son menos
satisfactorias en diversos aspectos que las de sus dos competidores (Exp.
XIV n\textsuperscript{o}~4). Así, suponiendo de nuevo que $f:X\to Y$ es un
morfismo propio de esquemas noetherianos regulares provistos de módulos
invertibles amplios, no es cierto, en general, que el homomorfismo
$f_*:K(X)\longrightarrow K(Y)$ aplique $\operatorname{Filt}^i K(X)$ en
$\operatorname{Filt}^{i-d} K(Y)$, donde $d$ es la dimensión relativa de $X$
sobre $Y$, y defina, por consiguiente, al pasar a los graduados asociados, un
homomorfismo de grado $-d$ de $\Gr^\bullet(X)$ en $\Gr^\bullet(Y)$. Sin
embargo, se puede demostrar (Exp. VII) que esto es así siempre que se
tensorice por $\mathbb Q$, de modo que se obtiene un homomorfismo de grado
$-d$
\begin{equation}\tag{2.5}
f_*:\Gr^\bullet(X)_{\mathbb Q}\longrightarrow \Gr^\bullet(Y)_{\mathbb Q},
\end{equation}
que desempeña el papel del homomorfismo «~imagen directa de ciclos~»
$A(X)\to A(Y)$.

Sentado esto, vemos que disponemos de todos los elementos de estructura
necesarios para dar sentido a la fórmula (1.2) cuando no se postula la
existencia de un cuerpo base, siendo $f:X\to Y$ un morfismo propio de
esquemas noetherianos regulares que admiten módulos invertibles amplios.
